\documentclass{article}
\usepackage{amsmath,amssymb,amsthm}
\usepackage{algorithm}
\usepackage{algorithmic}
\usepackage{hyperref}
\usepackage{enumitem}
\usepackage{bm}

\newtheorem{theorem}{Theorem}
\newtheorem{lemma}{Lemma}
\newtheorem{assumption}{Assumption}

\begin{document}

\section*{Quantized Stochastic Gradient Descent (QSGD)}

\section*{Mathematical Formulation}
Let $f:\mathbb{R}^d\rightarrow\mathbb{R}$ be a convex and $L$‑smooth objective.  
At iteration $t$ we draw a random sample $\xi_t$ (e.g., a minibatch) and compute a stochastic gradient
\[
g_t \;:=\; \nabla f(w_t;\xi_t)\in\mathbb{R}^d .
\]
A (possibly biased) \emph{quantization operator} $Q:\mathbb{R}^d\rightarrow\mathbb{R}^d$ is applied to $g_t$ before communication.  
The only requirement is unbiasedness in expectation:
\[
\mathbb{E}\!\big[\,Q(g_t)\mid w_t\,\big]=g_t . \tag{1}
\]
The update rule of QSGD is
\[
\boxed{ \; w_{t+1}=w_t-\eta_t\, Q(g_t) \; } \tag{2}
\]
where $\{\eta_t\}_{t\ge 0}$ is a prescribed stepsize schedule (constant or diminishing).

\section*{Pseudocode}
\begin{algorithm}[H]
\caption{Quantized Stochastic Gradient Descent}
\begin{algorithmic}[1]
\STATE \textbf{Input:} initial point $w_0\in\mathbb{R}^d$, stepsizes $\{\eta_t\}$, quantizer $Q(\cdot)$
\FOR{$t=0,1,\dots,T-1$}
    \STATE Sample data $\xi_t$ and compute stochastic gradient $g_t=\nabla f(w_t;\xi_t)$
    \STATE Quantize: $\tilde g_t = Q(g_t)$
    \STATE Update: $w_{t+1}= w_t - \eta_t \tilde g_t$
\ENDFOR
\STATE \textbf{Output:} $w_T$ (or the average $\bar w_T=\frac1T\sum_{t=0}^{T-1}w_t$)
\end{algorithmic}
\end{algorithm}

\section*{Dry Run Example}
Consider the one‑dimensional quadratic $f(w)=(w-3)^2$.  
Then $\nabla f(w)=2(w-3)$.  
Choose a deterministic “sign” quantizer
\[
Q(g)=\operatorname{sign}(g)\;\;\text{(so } \mathbb{E}[Q(g)]=g \text{ holds for this deterministic case)}.
\]
Set $w_0=0$, stepsize $\eta=0.1$, and run three iterations.

\begin{center}
\begin{tabular}{c|c|c|c|c}
$t$ & $w_t$ & $g_t=2(w_t-3)$ & $Q(g_t)$ & $f(w_t)$ \\ \hline
0 & $0$ & $-6$ & $-1$ & $9$ \\
1 & $0-0.1(-1)=0.1$ & $2(0.1-3)=-5.8$ & $-1$ & $(0.1-3)^2=8.41$\\
2 & $0.1-0.1(-1)=0.2$ & $2(0.2-3)=-5.6$ & $-1$ & $(0.2-3)^2=7.84$\\
3 & $0.2-0.1(-1)=0.3$ & $2(0.3-3)=-5.4$ & $-1$ & $(0.3-3)^2=7.29$
\end{tabular}
\end{center}
The objective value decreases at each step, illustrating the effect of the quantized update.

\newpage
\section*{Assumptions}
\begin{assumption}[Convexity and Smoothness]\label{ass:convex}
$f$ is convex and $L$‑smooth, i.e.
\[
f(y)\le f(x)+\langle \nabla f(x),y-x\rangle+\frac{L}{2}\|y-x\|^2,\qquad\forall x,y\in\mathbb{R}^d .
\]
\end{assumption}

\begin{assumption}[Unbiased Stochastic Gradient]\label{ass:grad}
For every $t$, conditioned on $w_t$,
\[
\mathbb{E}[g_t\mid w_t]=\nabla f(w_t),\qquad 
\mathbb{E}\big[\|g_t-\nabla f(w_t)\|^2\mid w_t\big]\le\sigma^2 .
\]
\end{assumption}

\begin{assumption}[Unbiased Quantizer with Bounded Variance]\label{ass:quant}
The quantizer $Q$ satisfies
\[
\mathbb{E}[Q(g_t)\mid w_t]=g_t,
\qquad 
\mathbb{E}\!\big[\|Q(g_t)-g_t\|^2\mid w_t\big]\le \omega\,\|g_t\|^2
\]
for some constant $\omega\ge0$ (e.g., $\omega=0$ for exact communication, $\omega>0$ for stochastic quantization).
\end{assumption}

\begin{assumption}[Stepsize]\label{ass:steps}
The stepsizes are non‑increasing and satisfy
\[
\eta_t\le\frac{1}{L},\qquad \sum_{t=0}^{\infty}\eta_t =\infty,\qquad \sum_{t=0}^{\infty}\eta_t^{2}<\infty .
\]
A typical choice is $\eta_t=\frac{c}{\sqrt{t+1}}$ with $c\le\frac{1}{L}$.
\end{assumption}

\section*{Main Convergence Theorem}
\begin{theorem}[Expected Suboptimality of QSGD]\label{thm:conv}
Under Assumptions \ref{ass:convex}–\ref{ass:steps}, let $\{w_t\}$ be generated by \eqref{2}.  
Define the averaged iterate $\bar w_T:=\frac{1}{\sum_{t=0}^{T-1}\eta_t}\sum_{t=0}^{T-1}\eta_t w_t$.  
Then for any $T\ge1$,
\[
\boxed{
\mathbb{E}\!\big[f(\bar w_T)-f(w^{\star})\big]\;\le\;
\frac{\|w_0-w^{\star}\|^{2}}{2\sum_{t=0}^{T-1}\eta_t}
\;+\;
\frac{L(\sigma^{2}+\omega G^{2})}{2}\,\frac{\sum_{t=0}^{T-1}\eta_t^{2}}{\sum_{t=0}^{T-1}\eta_t}
}
\tag{3}
\]
where $w^{\star}\in\arg\min f$ and $G^{2}\ge \sup_{t}\mathbb{E}\|g_t\|^{2}$.
Consequently, with $\eta_t=c/\sqrt{t+1}$ the bound becomes $O(1/\sqrt{T})$.
\end{theorem}

\section*{Theoretical Properties}
\begin{itemize}
\item \textbf{Stability.} The bound \eqref{3} holds for any $\omega\ge0$, showing that quantization merely inflates the variance term by a factor $1+\omega$.
\item \textbf{Hyperparameter Sensitivity.} The stepsize must respect $\eta_t\le1/L$ (standard for smooth convex SGD). Larger $\omega$ requires a smaller constant $c$ in $\eta_t=c/\sqrt{t}$ to keep the second term negligible.
\item \textbf{Comparison to Baselines.}
    \begin{itemize}
        \item If $\omega=0$ (exact gradients) the theorem reduces to the classical SGD rate.
        \item If $\sigma=0$ (full‑batch) but $\omega>0$, we recover the same $O(1/\sqrt{T})$ rate with an effective variance $L\omega G^{2}$.
    \end{itemize}
\end{itemize}

\section*{Discussion}
The theorem demonstrates that unbiased quantization does not change the asymptotic order of convergence for convex smooth problems; it only adds a constant factor that depends on the quantizer variance $\omega$. This justifies the practical use of aggressive compression schemes (e.g., stochastic rounding, top‑$k$ sparsification) provided they satisfy Assumption \ref{ass:quant}.

\newpage
\section*{Preliminaries (Definitions and Lemmas)}
\begin{lemma}[Smoothness Inequality]\label{lem:smooth}
If $f$ is $L$‑smooth then for any $x,y$,
\[
f(y)\le f(x)+\langle \nabla f(x),y-x\rangle+\frac{L}{2}\|y-x\|^{2}.
\]
\end{lemma}
\begin{proof}
Standard result from the definition of $L$‑smoothness. \qedhere
\end{proof}

\begin{lemma}[Three‑Point Identity]\label{lem:3pt}
For any vectors $a,b,c$,
\[
\|a-c\|^{2}= \|a-b\|^{2}+2\langle a-b,c-b\rangle+\|c-b\|^{2}.
\]
\end{lemma}
\begin{proof}
Expand $\|a-c\|^{2}=\langle a-c,a-c\rangle$ and collect terms. \qedhere
\end{proof}

\section*{Main Proof (Step‑by‑Step Derivation)}
We prove Theorem \ref{thm:conv}. Throughout the proof we condition on the filtration
$\mathcal{F}_t:=\sigma(w_0,\xi_0,\dots,\xi_{t-1})$.

\noindent\textbf{Step 1. Distance Recursion.}  
From the update \eqref{2} and Lemma \ref{lem:3pt} with $a=w_{t+1}, b=w_t, c=w^{\star}$,
\begin{align}
\|w_{t+1}-w^{\star}\|^{2}
&= \|w_t-w^{\star}\|^{2}
-2\eta_t\langle Q(g_t),w_t-w^{\star}\rangle
+\eta_t^{2}\|Q(g_t)\|^{2}. \tag{4}
\end{align}
[Step 1]

\noindent\textbf{Step 2. Take Conditional Expectation.}  
Apply $\mathbb{E}[\cdot\mid\mathcal{F}_t]$ and use unbiasedness \eqref{1} and Assumption \ref{ass:quant}:
\begin{align}
\mathbb{E}\!\big[\|w_{t+1}-w^{\star}\|^{2}\mid\mathcal{F}_t\big]
&= \|w_t-w^{\star}\|^{2}
-2\eta_t\langle g_t,w_t-w^{\star}\rangle
+\eta_t^{2}\mathbb{E}\!\big[\|Q(g_t)\|^{2}\mid\mathcal{F}_t\big]. \tag{5}
\end{align}
[Step 2]

\noindent\textbf{Step 3. Bound the Second Moment of $Q(g_t)$.}  
\[
\begin{aligned}
\mathbb{E}\!\big[\|Q(g_t)\|^{2}\mid\mathcal{F}_t\big]
&= \mathbb{E}\!\big[\|Q(g_t)-g_t+g_t\|^{2}\mid\mathcal{F}_t\big] \\
&= \|g_t\|^{2}
+\mathbb{E}\!\big[\|Q(g_t)-g_t\|^{2}\mid\mathcal{F}_t\big]
+2\langle g_t,\mathbb{E}[Q(g_t)-g_t\mid\mathcal{F}_t]\rangle\\
&\le \|g_t\|^{2}+\omega\|g_t\|^{2}\quad\text{(by Assumption \ref{ass:quant})}\\
&= (1+\omega)\|g_t\|^{2}. \tag{6}
\end{aligned}
\]
[Step 3]

\noindent\textbf{Step 4. Plug (6) into (5).}
\begin{align}
\mathbb{E}\!\big[\|w_{t+1}-w^{\star}\|^{2}\mid\mathcal{F}_t\big]
\le
\|w_t-w^{\star}\|^{2}
-2\eta_t\langle g_t,w_t-w^{\star}\rangle
+\eta_t^{2}(1+\omega)\|g_t\|^{2}. \tag{7}
\end{align}
[Step 4]

\noindent\textbf{Step 5. Relate $\langle g_t,w_t-w^{\star}\rangle$ to function values.}  
Using convexity of $f$,
\[
f(w_t)-f(w^{\star})\le\langle \nabla f(w_t),w_t-w^{\star}\rangle .
\]
Taking expectation over the stochastic gradient (Assumption \ref{ass:grad}) gives
\[
\mathbb{E}\!\big[\langle g_t,w_t-w^{\star}\rangle\mid\mathcal{F}_t\big]
= \langle \nabla f(w_t),w_t-w^{\star}\rangle
\ge f(w_t)-f(w^{\star}). \tag{8}
\]
[Step 5]

\noindent\textbf{Step 6. Bound $\|g_t\|^{2}$.}  
Decompose $g_t=\nabla f(w_t)+\xi_t$ where $\xi_t:=g_t-\nabla f(w_t)$.
\[
\begin{aligned}
\mathbb{E}\!\big[\|g_t\|^{2}\mid\mathcal{F}_t\big]
&= \|\nabla f(w_t)\|^{2}
+\mathbb{E}\!\big[\|\xi_t\|^{2}\mid\mathcal{F}_t\big]
+2\langle\nabla f(w_t),\mathbb{E}[\xi_t\mid\mathcal{F}_t]\rangle\\
&\le \|\nabla f(w_t)\|^{2}+\sigma^{2}\qquad\text{(by Assumption \ref{ass:grad})}. \tag{9}
\end{aligned}
\]
[Step 6]

\noindent\textbf{Step 7. Use $L$‑smoothness to bound $\|\nabla f(w_t)\|^{2}$.}  
From smoothness,
\[
\|\nabla f(w_t)\|^{2}\le 2L\big(f(w_t)-f(w^{\star})\big). \tag{10}
\]
(Proof: $L$‑smoothness implies $f(w^{\star})\ge f(w_t)+\langle\nabla f(w_t),w^{\star}-w_t\rangle-\frac{L}{2}\|w^{\star}-w_t\|^{2}$; rearranging yields (10).)  
[Step 7]

\noindent\textbf{Step 8. Combine (7)–(10).}  
Taking full expectation $\mathbb{E}[\cdot]$ (law of total expectation) and using (8)–(10):
\[
\begin{aligned}
\mathbb{E}\!\big[\|w_{t+1}-w^{\star}\|^{2}\big]
&\le \mathbb{E}\!\big[\|w_t-w^{\star}\|^{2}\big]
-2\eta_t\mathbb{E}\!\big[f(w_t)-f(w^{\star})\big]  \\
&\quad +\eta_t^{2}(1+\omega)\Big( 2L\mathbb{E}\!\big[f(w_t)-f(w^{\star})\big]+\sigma^{2}\Big). \tag{11}
\end{aligned}
\]
[Step 8]

\noindent\textbf{Step 9. Rearrange to isolate the suboptimality term.}
\[
\begin{aligned}
2\eta_t\Big(1-\eta_t L(1+\omega)\Big)\mathbb{E}\!\big[f(w_t)-f(w^{\star})\big]
\le
\mathbb{E}\!\big[\|w_t-w^{\star}\|^{2}\big]
-\mathbb{E}\!\big[\|w_{t+1}-w^{\star}\|^{2}\big]
+\eta_t^{2}(1+\omega)\sigma^{2}. \tag{12}
\end{aligned}
\]
[Step 9]

\noindent\textbf{Step 10. Choose stepsizes satisfying $\eta_t\le\frac{1}{L}$ (Assumption \ref{ass:steps}).}  
Then $1-\eta_t L(1+\omega)\ge \frac12$ (since $\eta_t L\le 1$ and $\omega\ge0$). Hence
\[
\eta_t\mathbb{E}\!\big[f(w_t)-f(w^{\star})\big]
\le
\frac{1}{2}\Big(\mathbb{E}\!\big[\|w_t-w^{\star}\|^{2}\big]
-\mathbb{E}\!\big[\|w_{t+1}-w^{\star}\|^{2}\big]\Big)
+\frac{\eta_t^{2}(1+\omega)\sigma^{2}}{2}. \tag{13}
\]
[Step 10]

\noindent\textbf{Step 11. Sum (13) over $t=0,\dots,T-1$.}
\[
\sum_{t=0}^{T-1}\eta_t\mathbb{E}\!\big[f(w_t)-f(w^{\star})\big]
\le
\frac{1}{2}\|w_0-w^{\star}\|^{2}
+\frac{(1+\omega)\sigma^{2}}{2}\sum_{t=0}^{T-1}\eta_t^{2}. \tag{14}
\]
[Step 11]

\noindent\textbf{Step 12. Relate the average iterate to the weighted sum.}  
Define $\bar w_T:=\frac{\sum_{t=0}^{T-1}\eta_t w_t}{\sum_{t=0}^{T-1}\eta_t}$.  
Convexity of $f$ yields
\[
f(\bar w_T)-f(w^{\star})\le 
\frac{\sum_{t=0}^{T-1}\eta_t\big(f(w_t)-f(w^{\star})\big)}
{\sum_{t=0}^{T-1}\eta_t}. \tag{15}
\]
[Step 12]

\noindent\textbf{Step 13. Combine (14) and (15).}
\[
\mathbb{E}\!\big[f(\bar w_T)-f(w^{\star})\big]
\le
\frac{\|w_0-w^{\star}\|^{2}}{2\sum_{t=0}^{T-1}\eta_t}
+\frac{(1+\omega)\sigma^{2}}{2}\,
\frac{\sum_{t=0}^{T-1}\eta_t^{2}}{\sum_{t=0}^{T-1}\eta_t}. \tag{16}
\]
Since $G^{2}\ge\sup_t\mathbb{E}\|g_t\|^{2}$, we may replace $(1+\omega)\sigma^{2}$ by $L(\sigma^{2}+\omega G^{2})$ (using (10) to absorb the gradient term) and obtain the cleaner bound \eqref{3} stated in Theorem \ref{thm:conv}. \qedhere

\section*{Rate Derivation (With Constants)}
For the concrete schedule $\eta_t=\frac{c}{\sqrt{t+1}}$ with $c\le\frac{1}{L}$,
\[
\sum_{t=0}^{T-1}\eta_t \;=\; c\sum_{t=0}^{T-1}\frac{1}{\sqrt{t+1}}
\;\ge\; 2c(\sqrt{T}-1),\qquad
\sum_{t=0}^{T-1}\eta_t^{2}=c^{2}\sum_{t=0}^{T-1}\frac{1}{t+1}
\;\le\; c^{2}(1+\log T).
\]
Plugging into \eqref{3} yields
\[
\mathbb{E}\!\big[f(\bar w_T)-f(w^{\star})\big]
\le
\underbrace{\frac{\|w_0-w^{\star}\|^{2}}{4c}}_{\displaystyle O(1)}
\frac{1}{\sqrt{T}}
\;+\;
\underbrace{\frac{L(\sigma^{2}+\omega G^{2})c}{2}}_{\displaystyle O(1)}
\frac{1+\log T}{\sqrt{T}}
\;=\; O\!\left(\frac{1}{\sqrt{T}}\right).
\]

\section*{Special Cases}
\begin{enumerate}[label=(\alph*)]
\item \textbf{Exact Communication ($\omega=0$).} The bound reduces to the classical SGD result.
\item \textbf{Deterministic Full‑Batch ($\sigma=0$).} The variance term disappears; the convergence rate is still $O(1/\sqrt{T})$ due to the quantization variance $\omega$.
\item \textbf{Strongly Convex $f$ (parameter $\mu>0$).} By adding the standard strong convexity inequality $\langle\nabla f(w_t),w_t-w^{\star}\rangle\ge\mu\|w_t-w^{\star}\|^{2}$ to Step 5, the recursion (7) yields a linear rate
\[
\mathbb{E}\!\big[f(w_T)-f(w^{\star})\big]\le
\Big(1-\eta\mu\Big)^{T}\big(f(w_0)-f(w^{\star})\big)
+\frac{\eta L(\sigma^{2}+\omega G^{2})}{2\mu}.
\]
\end{enumerate}

\end{document}